% !TEX root = ../main.tex
\documentclass[../main.tex]{subfiles}

\begin{document}

\chapter{Software and Reproducibility}
\labch{appendix-software}

This appendix documents the software tools, code repositories, and reproducibility materials associated with this dissertation. All code is available under the \RoleBox\ umbrella repository.

\section{Repository Structure}

The dissertation materials are organized as a collection of Git repositories with the main \RoleBox\ repository serving as an umbrella:

\begin{kaobox}[frametitle=Repository Architecture]
\begin{verbatim}
rolebox/                     # Main umbrella repository
├── dissertation/            # LaTeX source (kaobook)
└── projects/                # Analysis repositories
    ├── rolesim/             # Agent-based simulation
    ├── rolexray/            # Role extraction & NLP
    ├── rolebox-wikipedia/   # Wikipedia/Wikidata data
    ├── rolebox-websites/    # Website corpus
    ├── rolebox-twitter/     # Twitter/X data
    ├── rolebox-images/      # Visual corpus
    ├── rolebox-crunchbase/  # Investment data
    ├── rolebox-news/        # News media corpus
    └── rolebox-video/       # Video content
\end{verbatim}
\end{kaobox}

\section{Core Tools}

\begin{description}[leftmargin=0.5cm]

\item[\RoleBox] Main umbrella repository containing all dissertation materials\\
\faGithub\ \texttt{github.com/babajideowoyele/rolebox}

\item[\RoleSim] Agent-based role dynamics simulation framework (Python)\\
\faGithub\ \texttt{github.com/babajideowoyele/rolesim}\\
Features: Influence typology, adaptive agent strategies, network metrics

\item[RoleXray] Role extraction and classification toolkit (Python)\\
\faGithub\ \texttt{github.com/babajideowoyele/rolexray}\\
Features: BERTopic integration, spaCy NLP, discourse network analysis

\end{description}

\section{Data Repositories}

Each data modality has a dedicated repository with collection scripts, processed data, and interactive exploration UIs:

\begin{description}[leftmargin=0.5cm]

\item[rolebox-wikipedia] Wikipedia and Wikidata entity collection\\
\faGithub\ \texttt{github.com/babajideowoyele/rolebox-wikipedia}\\
Content: 112 NetZeroCities, 133 EIT KIC partner organizations

\item[rolebox-websites] KIC partner website corpus\\
\faGithub\ \texttt{github.com/babajideowoyele/rolebox-websites}

\item[rolebox-twitter] Twitter/X social media data\\
\faGithub\ \texttt{github.com/babajideowoyele/rolebox-twitter}

\item[rolebox-images] Visual materials corpus\\
\faGithub\ \texttt{github.com/babajideowoyele/rolebox-images}

\item[rolebox-crunchbase] Venture and investment data\\
\faGithub\ \texttt{github.com/babajideowoyele/rolebox-crunchbase}

\item[rolebox-news] News media corpus\\
\faGithub\ \texttt{github.com/babajideowoyele/rolebox-news}

\item[rolebox-video] Video content analysis\\
\faGithub\ \texttt{github.com/babajideowoyele/rolebox-video}

\end{description}

\section{Software Dependencies}

The toolkit is built on the following major dependencies:

\begin{itemize}
\item \textbf{Python 3.11+} -- Primary programming language
\item \textbf{BERTopic} -- Topic modeling and document clustering
\item \textbf{sentence-transformers} -- Text embeddings
\item \textbf{NetworkX} -- Network analysis
\item \textbf{Plotly} -- Interactive visualizations
\item \textbf{Pandas/NumPy} -- Data manipulation
\item \textbf{Carbon Design System} -- UI components
\end{itemize}

\section{Interactive UIs}

Each data repository includes an interactive HTML/JavaScript interface for exploring the data:

\begin{itemize}
\item Built with IBM Carbon Design System for consistent styling
\item Plotly.js for interactive charts and maps
\item Sigma.js for network visualizations
\item All UIs can be run locally via \texttt{python -m http.server}
\end{itemize}

\section{Reproducibility}

To reproduce the analyses:

\begin{enumerate}
\item Clone the umbrella repository with submodules:\\
\texttt{git clone --recurse-submodules github.com/babajideowoyele/rolebox}

\item Create conda environments for each project:\\
\texttt{conda env create -f environment.yml}

\item Follow the notebooks in each project's \texttt{notebooks/} folder
\end{enumerate}

\section{Data Availability}

Raw data is available in each repository's \texttt{data/} folder, subject to the following:

\begin{itemize}
\item Wikipedia/Wikidata: CC BY-SA 4.0
\item Twitter data: Subject to Twitter Terms of Service
\item Crunchbase: Subject to Crunchbase licensing
\item News media: Subject to publisher agreements
\end{itemize}

\end{document}
